% paper2_methods.tex -- Materials and Methods of Paper 2, input by
% papers/paper2_nitrogen.tex.
\section{Materials and Methods}

\subsection{Treatments and design}
Maize received five rates of mineral nitrogen, 0, 40, 80, 120 and
\SI{160}{\kilo\gram\per\hectare}, applied as urea in two equal splits at
sowing and at the six-leaf stage. Phosphorus and potassium were applied to all
plots at rates that did not limit growth. The rates were arranged in
randomised complete blocks with \res{design-blocks} blocks at each of the three sites
\parencite{gomez1984}. Grain yield is expressed in \si{\tonne\per\hectare} on a
dry-matter basis, and nitrogen uptake as the nitrogen in grain and stover at
harvest. The data are held in \texttt{data/trials/t2\_nitrogen.tsv}.

\subsection{Yield response}
The response of grain yield $Y$ to the nitrogen rate $N$ was described by a
quadratic-plateau\index{quadratic-plateau response} function,
\begin{equation}
  Y =
  \begin{cases}
    a + bN + cN^2 & N < N_p,\\
    Y_p & N \ge N_p,
  \end{cases}
  \qquad N_p = -\frac{b}{2c},
  \label{eq:quadplateau}
\end{equation}
where $a$, $b$ and $c$ are fitted coefficients ($c < 0$), $N_p$ is the rate at
which the response reaches its plateau and $Y_p$ is the plateau yield. $N_p$ is
taken as the agronomic optimum rate\index{nitrogen!agronomic optimum rate}.

\subsection{Nitrogen use efficiency}
Two efficiency measures were calculated for each fertilised plot against the
unfertilised plot of the same block \parencite{cassman2002,ladha2005}. The
agronomic efficiency of applied nitrogen, in kilograms of grain per kilogram of
nitrogen, is
\begin{equation}
  \mathrm{AE}_N = \frac{1000\,(Y_N - Y_0)}{N},
  \label{eq:ae}
\end{equation}
where $Y_N$ and $Y_0$ are the grain yields (\si{\tonne\per\hectare}) with and
without nitrogen. The apparent recovery\index{nitrogen use efficiency!recovery efficiency} efficiency, the fraction of the applied
nitrogen taken up by the crop, is
\begin{equation}
  \mathrm{RE}_N = \frac{U_N - U_0}{N},
  \label{eq:re}
\end{equation}
where $U_N$ and $U_0$ are the nitrogen uptakes (\si{\kilo\gram\per\hectare})
with and without nitrogen.
